% GENERATED FILE. Edit canonical lessons, then run npm run generate:books.
% canonical-source-hash: fnv1a64:282be41f6584f28e
% canonical-lessons: ZH-C392-gai, ZH-W392-gai-alter, ZH-C392-gaiqi, ZH-C392-gaibian, ZH-C392-gaiming, ZH-C392-gaikou

\chapter{A Change of Plan}
\label{ch:zh-a-change-of-plan}
\hlchaptermodality{392}

\begin{chapteropening}
\textbf{By the end of this chapter:} \emph{I can talk about changing things, changing a date, changing plans, new names and someone changing their tune.}

Complete the last lesson of chapter 392: I can talk about changing things, changing a date, changing plans, new names and someone changing their tune.
\end{chapteropening}

\section[\texorpdfstring{\zh{改} (gǎi)}{改 (gǎi)}]{\zh{改} (gǎi) — to alter; to amend}

\label{lesson:ZH-C392-gai}



\begin{quote}
Before the new one: say the Mandarin for to make it onto a train, then the Mandarin for when the time comes.
\end{quote}

\subsection*{You'll want to know: \zh{改}}
\textbf{\zh{改}} — \emph{gǎi} — \textquotedblleft{}to alter; to amend\textquotedblright{}.

\begin{quote}
\textbf{\zh{时间改了我们下午三点见面}} \emph{shíjiān gǎi le wǒmen xiàwǔ sān diǎn jiànmiàn} — The time's changed: we're meeting at three in the afternoon.
\end{quote}

The new character, in pieces: \zh{改} — seven strokes in two pieces: on the left a bar that turns down, a short bar, and a vertical that bends right and flicks up; on the right the right piece of \zh{做}. The next lesson puts it in your hand.

\begin{cousinweb}
\textbf{\zh{改}} (\emph{gǎi}) is inherited from Middle Chinese \emph{kojX} (Guangyun, in Baxter's transcription), glossed 'to change'.
\end{cousinweb}

\begin{grammarlens}[title={What you've built}]
One more word for saying what would happen.
\end{grammarlens}

\subsection*{Your turn}
\begin{itemize}
\raggedright
  \item \emph{Say it:} \emph{gǎi}
  \item \emph{Say it:} \emph{gǎi}, once more
  \item \emph{Say it:} \emph{shíjiān gǎi le wǒmen xiàwǔ sān diǎn jiànmiàn} — The time's changed: we're meeting at three in the afternoon.
  \item \emph{Recall:} say the Mandarin for to make it onto a train, then the Mandarin for when the time comes, then say \emph{gǎi} again
\end{itemize}

\begin{tcolorbox}[breakable,colback=teal!4,colframe=teal!35!black,title={Before you move on}]
What does \zh{改} mean? (\textquotedblleft{}To alter; to amend\textquotedblright{}.) Say it once more.
\end{tcolorbox}
\section[\texorpdfstring{\zh{改} (gǎi)}{改 (gǎi)}]{\zh{改} — a bent piece, then the right piece of \zh{做}}

\label{lesson:ZH-W392-gai-alter}



\begin{quote}
The word first, then the character inside it.

\begin{itemize}
\raggedright
  \item \emph{Recall:} say \emph{gǎi} — \textbf{R1}, one lesson back, before it decays
\end{itemize}

\begin{quote}
\zh{做} \zh{变}
\end{quote}

\emph{Read it:} each one aloud, with nothing said first — \emph{zuò}, \emph{biàn}
\end{quote}

\subsection*{Script — a bent piece, then the right piece of \zh{做}}
\hlblockfigure{\detokenize{figures/ZH-W392-gai-alter-filmstrip.pdf}}{How \zh{改} is written, stroke by stroke}

\begin{quote}
\zh{改}
\end{quote}

\textbf{Two pieces, left before right.}

\textbf{The left piece, three strokes.} A bar that turns down. A short bar under it. Then a vertical that comes down, bends right along the foot and flicks up.

\textbf{The right piece, four strokes, the right piece of \zh{做}.} A short stroke falling to the left. A bar from it to the right. A long stroke falling to the left. Then a long stroke falling to the right across it.

3 + 4 = seven.

\subsection*{Your turn}
\begin{itemize}
\raggedright
  \item \emph{Write it:} \zh{改}, all seven, with the model in front of you
  \item \emph{Cover:} write it again from memory, and count the pieces, not the strokes
\end{itemize}

\begin{tcolorbox}[breakable,colback=teal!4,colframe=teal!35!black,title={Before you move on}]
How many strokes in \zh{改}? (\textbf{Seven}.) Where have you made the right piece before? (\textbf{On the right of \zh{做}}.) How does the left piece end? (\textbf{A bend to the right and a flick up}.)

Source: \href{https://raw.githubusercontent.com/chanind/hanzi-writer-data/68d10a4b21150cae5e1ebbd223eed289cf32d90c/data/\%E6\%94\%B9.json}{Hanzi Writer Data \zh{改}, pinned stroke paths and medians}.
\end{tcolorbox}
\section[\texorpdfstring{\zh{改期} (gǎiqī)}{改期 (gǎiqī)}]{\zh{改期} (gǎiqī) — to change the date (of an event)}

\label{lesson:ZH-C392-gaiqi}



\begin{quote}
Before the new one: say the Mandarin for when the time comes, then the Mandarin for to alter.
\end{quote}

\subsection*{You'll want to know: \zh{改期}}
\textbf{\zh{改期}} — \emph{gǎiqī} — \textquotedblleft{}to change the date (of an event)\textquotedblright{}.

\begin{quote}
\textbf{\zh{要是下雨晚会就改期}} \emph{yàoshi xiàyǔ wǎnhuì jiù gǎiqī} — If it rains, the party will be moved to another date.
\end{quote}

\begin{cousinweb}
\textbf{\zh{改期}} (\emph{gǎiqī}) is built from \zh{改} 'to alter' and \zh{期} 'a set time', both inherited from Middle Chinese, \emph{kojX} and \emph{gi} (Guangyun, in Baxter's transcription), which glosses the second 'an appointed meeting'.
\end{cousinweb}

\begin{grammarlens}[title={What you've built}]
One more word for saying what would happen.
\end{grammarlens}

\subsection*{Your turn}
\begin{itemize}
\raggedright
  \item \emph{Say it:} \emph{gǎiqī}
  \item \emph{Say it:} \emph{gǎiqī}, once more
  \item \emph{Say it:} \emph{yàoshi xiàyǔ wǎnhuì jiù gǎiqī} — If it rains, the party will be moved to another date.
  \item \emph{Recall:} say the Mandarin for when the time comes, then the Mandarin for to alter, then say \emph{gǎiqī} again
\end{itemize}

\begin{tcolorbox}[breakable,colback=teal!4,colframe=teal!35!black,title={Before you move on}]
What does \zh{改期} mean? (\textquotedblleft{}To change the date (of an event)\textquotedblright{}.) Say it once more.
\end{tcolorbox}
\section[\texorpdfstring{\zh{改变} (gǎibiàn)}{改变 (gǎibiàn)}]{\zh{改变} (gǎibiàn) — to make different (a plan or a habit)}

\label{lesson:ZH-C392-gaibian}



\begin{quote}
Before the new one: say the Mandarin for to alter, then the Mandarin for to change the date.
\end{quote}

\subsection*{You'll want to know: \zh{改变}}
\textbf{\zh{改变}} — \emph{gǎibiàn} — \textquotedblleft{}to make different (a plan or a habit)\textquotedblright{}.

\begin{quote}
\emph{rúguǒ tiānqì bù hǎo wǒmen jiù gǎibiàn jìhuà} — If the weather's bad, we'll change our plans.
\end{quote}

\begin{cousinweb}
\textbf{\zh{改变}} (\emph{gǎibiàn}) is built from \zh{改} 'to alter' and \zh{变} 'to change', both inherited from Middle Chinese, \emph{kojX} and \emph{pjenH} (Guangyun, in Baxter's transcription), which glosses the first 'to change' and the second 'to transform'.
\end{cousinweb}

\begin{grammarlens}[title={What you've built}]
One more word for saying what would happen.
\end{grammarlens}

\subsection*{Your turn}
\begin{itemize}
\raggedright
  \item \emph{Say it:} \emph{gǎibiàn}
  \item \emph{Say it:} \emph{gǎibiàn}, once more
  \item \emph{Say it:} \emph{rúguǒ tiānqì bù hǎo wǒmen jiù gǎibiàn jìhuà} — If the weather's bad, we'll change our plans.
  \item \emph{Recall:} say the Mandarin for to alter, then the Mandarin for to change the date, then say \emph{gǎibiàn} again
\end{itemize}

\begin{tcolorbox}[breakable,colback=teal!4,colframe=teal!35!black,title={Before you move on}]
What does \zh{改变} mean? (\textquotedblleft{}To make different (a plan or a habit)\textquotedblright{}.) Say it once more.
\end{tcolorbox}
\section[\texorpdfstring{\zh{改名} (gǎimíng)}{改名 (gǎimíng)}]{\zh{改名} (gǎimíng) — to change the name (of)}

\label{lesson:ZH-C392-gaiming}



\begin{quote}
Before the new one: say the Mandarin for to change the date, then the Mandarin for to make different.
\end{quote}

\subsection*{You'll want to know: \zh{改名}}
\textbf{\zh{改名}} — \emph{gǎimíng} — \textquotedblleft{}to change the name (of)\textquotedblright{}.

\begin{quote}
\textbf{\zh{那家饭店改名了}} \emph{nà jiā fàndiàn gǎimíng le} — That restaurant has changed its name.
\end{quote}

\begin{cousinweb}
\textbf{\zh{改名}} (\emph{gǎimíng}) is built from \zh{改} 'to alter' and \zh{名} 'a name', both inherited from Middle Chinese, \emph{kojX} and \emph{mjieng} (Guangyun, in Baxter's transcription).
\end{cousinweb}

\begin{grammarlens}[title={What you've built}]
One more word for saying what would happen.
\end{grammarlens}

\subsection*{Your turn}
\begin{itemize}
\raggedright
  \item \emph{Say it:} \emph{gǎimíng}
  \item \emph{Say it:} \emph{gǎimíng}, once more
  \item \emph{Say it:} \emph{nà jiā fàndiàn gǎimíng le} — That restaurant has changed its name.
  \item \emph{Recall:} say the Mandarin for to change the date, then the Mandarin for to make different, then say \emph{gǎimíng} again
\end{itemize}

\begin{tcolorbox}[breakable,colback=teal!4,colframe=teal!35!black,title={Before you move on}]
What does \zh{改名} mean? (\textquotedblleft{}To change the name (of)\textquotedblright{}.) Say it once more.
\end{tcolorbox}
\section[\texorpdfstring{\zh{改口} (gǎikǒu)}{改口 (gǎikǒu)}]{\zh{改口} (gǎikǒu) — to change one's tune; to take back what one said}

\label{lesson:ZH-C392-gaikou}



\begin{quote}
Before the new one: say the Mandarin for to make different, then the Mandarin for to change the name (of).
\end{quote}

\subsection*{You'll want to know: \zh{改口}}
\textbf{\zh{改口}} — \emph{gǎikǒu} — \textquotedblleft{}to change one's tune; to take back what one said\textquotedblright{}.

\begin{quote}
\emph{tā zuótiān shuō bú qù jīntiān yòu gǎikǒu le} — Yesterday he said he wasn't going; today he's changed his tune.
\end{quote}

Two third tones meet in this word: gǎi is said with a rising tone before kǒu, and the spelling keeps both marks.

\begin{cousinweb}
\textbf{\zh{改口}} (\emph{gǎikǒu}) is built from \zh{改} 'to alter' and \zh{口} 'the mouth', both inherited from Middle Chinese, \emph{kojX} and \emph{khuwX} (Guangyun, in Baxter's transcription), which glosses the second 'what a person speaks and eats with': altering what comes out of the mouth.
\end{cousinweb}

\begin{grammarlens}[title={What you've built}]
One more word for saying what would happen.
\end{grammarlens}

\subsection*{Your turn}
\begin{itemize}
\raggedright
  \item \emph{Say it:} \emph{gǎikǒu}
  \item \emph{Say it:} \emph{gǎikǒu}, once more
  \item \emph{Say it:} \emph{tā zuótiān shuō bú qù jīntiān yòu gǎikǒu le} — Yesterday he said he wasn't going; today he's changed his tune.
  \item \emph{Recall:} say the Mandarin for to make different, then the Mandarin for to change the name (of), then say \emph{gǎikǒu} again
\end{itemize}

\begin{tcolorbox}[breakable,colback=teal!4,colframe=teal!35!black,title={Before you move on}]
What does \zh{改口} mean? (\textquotedblleft{}To change one's tune; to take back what one said\textquotedblright{}.) Say it once more.
\end{tcolorbox}
